\chapter{Limites, perspectives, conclusion}
\label{ch:limites}

\section{Limites du travail}

\subsection{Étude de cas unique}

Tout ce qui est mesuré l'est sur \emph{un seul} environnement~:
le lab UrbanUpC. On n'a pas tourné l'outil sur un autre ENT, ni
sur un SI de taille différente. Les chiffres sont valables pour
ce lab~; leur transposition à un autre contexte demande au
minimum~:

\begin{itemize}
  \item un déploiement sur une autre infrastructure cloud
    (autre RG Azure, ou autre fournisseur)~;
  \item une adaptation des watchers à l'AD réel et au SIEM
    réel de l'organisation cible~;
  \item une nouvelle classification A/B/C des preuves attendues
    pour le périmètre métier concerné.
\end{itemize}

\subsection{Volume modeste}

Le lab fait tourner 3 VMs, 3 applications, une vingtaine
d'utilisateurs AD. Un ENT en production a plutôt~:

\begin{itemize}
  \item 50 à 200 applications homologuées~;
  \item 10\,000 à 30\,000 utilisateurs AD actifs~;
  \item plusieurs centaines de règles NSG/firewall.
\end{itemize}

Les 90~secondes par cycle deviendraient probablement 10 à
30~minutes à cette échelle. Rien d'impossible, mais il faudrait
paralléliser les watchers et probablement passer SQLite à
PostgreSQL pour le stockage.

\subsection{Couverture technique uniquement}

L'outil n'attaque que la couche technique. Les preuves
organisationnelles (formations effectuées, revues d'analyse de
risque, procédures appliquées) restent manuelles. Un dossier
HOMO-CI seul ne suffit pas à homologuer~: il alimente une part
du dossier, l'autre part vient du chef de projet et de l'AQSSI.

\subsection{Pas encore validé par l'ANSSI}

Le cadre actuel reconnaît un dossier daté et signé. Rien dans le
texte ne dit qu'un dossier régénéré chaque nuit est conforme. À
ce stade, HOMO-CI produit un livrable \emph{compatible avec la
structure ANSSI}, mais pas un dossier officiellement reconnu.
Une démarche de discussion avec l'ANSSI pour valider ce modèle
hybride serait un prolongement naturel.

\subsection{Watchers en stub}

En v0.1, deux watchers sont complets (\texttt{wc\_nsg} et
\texttt{wc\_cve}). Les watchers AD et Wazuh sont câblés dans
l'architecture mais retournent des données simulées. Les
compléter est la première chose à faire pour une v0.2.

\section{Perspectives}

\subsection{Court terme (3--6 mois)}

\begin{itemize}
  \item Finir les watchers \texttt{wc\_ad} (énumération LDAP des
    comptes, groupes, ACL) et \texttt{wc\_siem} (API Wazuh~:
    règles actives, état des agents, alertes par catégorie).
  \item Ajouter un watcher \texttt{wc\_asvs} pour cartographier
    automatiquement la conformité OWASP~ASVS sur les
    applications.
  \item Ajouter une boucle de remédiation simple~: pour les
    dérives à matérialité \emph{high} (NSG ouvert sur Internet,
    CVE critique), pouvoir proposer un PR git automatique.
\end{itemize}

\subsection{Moyen terme (6--18 mois)}

\begin{itemize}
  \item Validation sur un vrai ENT. L'idéal serait un partenariat
    avec une université qui accepte de déployer HOMO-CI sur un
    périmètre limité (1 ou 2 applications).
  \item Intégration avec un système de tickets (Jira,
    GitLab~Issues) pour que chaque dérive détectée crée
    automatiquement un ticket assigné au responsable.
  \item Interface web minimaliste pour permettre aux AQSSI non
    techniques de consulter l'état du dossier sans passer par
    le CLI.
\end{itemize}

\subsection{Long terme et transposition}

Le pipeline est conçu pour un ENT, mais la mécanique
(\emph{watchers, stockage, diff, rendu signé}) est
indépendante du domaine métier. Les domaines naturels de
transposition~:

\begin{itemize}
  \item \textbf{Transport.} Les SI des opérateurs de transport
    (RATP, SNCF) sont des cibles NIS~2 avec des contraintes
    similaires. Les watchers seraient à adapter (équipements
    industriels, supervision OT) mais le modèle reste valable.
  \item \textbf{Santé et hôpitaux.} Données sensibles, cadre
    légal contraignant, infra hétérogène --- terrain idéal pour
    de la conformité continue.
  \item \textbf{Collectivités territoriales.} Périmètre proche
    de celui d'un ENT, mêmes contraintes ANSSI/NIS~2.
\end{itemize}

\section{Ce qu'on retient}

\subsection{La technique}

Le pipeline est modeste en taille (\textasciitilde{}2\,000~lignes
de Python) et utilise des briques matures. Sa valeur n'est pas
dans une innovation algorithmique~: elle est dans
\textbf{l'assemblage}, et dans la rigueur du modèle de données
(\emph{Evidence} avec hash, chaînage SHA-256 entre dossiers,
canonicalisation JSON).

\subsection{Le lab}

UrbanUpC est, en lui-même, un livrable utile pour la formation~:
un SOC complet, reproductible, avec des attaques documentées et
des règles de détection qui fonctionnent. Il sert ici de banc de
test pour HOMO-CI, mais pourrait servir de support à des TPs
cybersécurité sans modification majeure.

\subsection{Le constat principal}

\emph{Maintenir un dossier d'homologation à jour est possible,
sans effort proportionnel à la taille du SI, à condition de
limiter le périmètre à la couche technique.} C'est le résultat
le plus important du PFE~: il ne défend pas l'automatisation
totale, qui n'a pas de sens (les preuves de gouvernance restent
humaines), mais il démontre que la partie automatisable peut
l'être avec un outil de taille raisonnable.

\subsection{Et ce qui reste à faire}

Beaucoup de choses, dont les plus importantes~: finir les
watchers manquants, valider sur un autre environnement,
discuter avec l'ANSSI pour valider la compatibilité du modèle,
écrire un guide d'adoption pour qu'un autre étudiant ou un autre
ingénieur puisse reprendre le projet en quelques jours.

\begin{takeaway}
Le PFE livre deux choses concrètes~: un lab SOC complet
(UrbanUpC, 3~VMs Azure, Wazuh, AD, applications vulnérables) et
un outil d'automatisation (HOMO-CI, pipeline Python en six
couches). L'outil fait son travail sur la couche technique~; il
ne prétend pas remplacer l'homologation officielle, juste la
rendre plus tenable dans le temps.
\end{takeaway}
